\section{Introducción}

\subsection{Motivación y Planteamiento del Problema}

Las arquitecturas transformer son el paradigma dominante para el modelado de secuencias en procesamiento del lenguaje natural \citep{vaswani_attention_2017}, visión por computadora y biología computacional, con modelos como BERT \citep{devlin_bert_2019} y GPT estableciendo el estándar para el aprendizaje de representaciones. Desde entonces, el campo se ha diversificado rápidamente: bloques recurrentes y retentivos, atención dispersa, mecanismos con compuerta, proyecciones latentes, capas con memoria, reglas de peso rápido y formulaciones de profundidad continua abordan cada uno una limitación distinta de la atención densa, mientras que el panorama de optimizadores se ha ampliado más allá de AdamW con métodos de reducción de varianza, eficientes en memoria, sin programación, de segundo orden y de bajo rango \citep{child_sparse_transformer_2019,sun_retentive_2023,gu_mamba_2023,behrouz_titans_2025,ramapuram_theory_2024,xie_adan_2022,pagliardini_ademamix_2024,yuan_mars_2024,shazeer_adafactor_2018,zhao_galore_2024,defazio_road_2024,gupta_shampoo_2018,vyas_soap_2024}.

La consecuencia práctica es un ecosistema fragmentado. Comparar mezcladores de secuencia u optimizadores bajo condiciones de entrenamiento idénticas exige reimplementar cada variante, integrarla en un pipeline de entrenamiento y mantener el resultado ante cambios del framework. Este esfuerzo dificulta la reproducibilidad, ralentiza el progreso científico y eleva la barrera de entrada para nuevos investigadores.

\subsection{Propuesta y Contribuciones}

Proponemos \textbf{Frankenstein Transformer} (\url{https://github.com/erickfmm/frankenstein-transformer}), un conjunto de herramientas basado en configuración que trata las decisiones de arquitectura, optimizador y entrenamiento como un único documento validado. El proyecto ofrece:

\begin{enumerate}[leftmargin=*]
    \item \textbf{Diseño basado en esquema}: un contrato YAML estricto que soporta cuarenta y tres mezcladores de secuencia en ocho familias y veinticuatro familias de optimizadores, validado antes de iniciar el entrenamiento.
    \item \textbf{Entrenamiento agnóstico a la arquitectura}: un único pipeline para líneas base densas, bloques recurrentes y retentivos, patrones dispersos, mecanismos con compuerta, compresión latente, memoria condicional, atención de peso rápido y atención de campo geométrico, con modos encoder y decoder \citep{zhu2026mixtureofdepthsattention,cheng2026conditionalmemoryscalablelookup,xie_mhc_2025,zhang_falcon_2026,pisoni_ssog_2026}.
    \item \textbf{Enrutamiento de optimizadores}: grupos de hiperparámetros con prefijo que controlan de forma independiente embeddings, capas de normalización, bloques de atención y el resto de parámetros para cada familia de optimizadores \citep{zhu2025apollosgdlikememoryadamwlevel}.
    \item \textbf{Flujos de trabajo de extremo a extremo}: entrenamiento (incluidos objetivos de embeddings de oraciones), despliegue y exportación (formatos cuantizado, estándar, HuggingFace Transformers y BitNet GGUF), inferencia y un constructor web generado desde el esquema, inspirado en flujos de trabajo estilo SBERT \citep{reimers_sentence-bert_2019}.
    \item \textbf{Soporte de visión}: una clase de modelo Vision Transformer que cubre predicción de parches enmascarados, clasificación y segmentación bajo la misma infraestructura de mezcladores y optimizadores.
    \item \textbf{Herramientas de validación}: pruebas unitarias, validación de presets YAML e integración continua en las versiones de Python soportadas.
\end{enumerate}

La superficie de comandos es \texttt{frankenstein-transformer}, con cuatro subcomandos: \texttt{train}, \texttt{deploy}, \texttt{infer} y \texttt{web-server}. El último lanza un constructor de configuración en Streamlit que renderiza los metadatos del esquema como formularios, valida la entrada en vivo, previsualiza el YAML resultante y genera el comando CLI equivalente, sirviendo tanto de ayuda a la configuración como de documentación del espacio de diseño disponible.

\subsection{Guía de Lectura}

Este documento se organiza como una referencia técnica:
\begin{enumerate}[leftmargin=*]
    \item \textbf{Requisitos previos}: el Apéndice~\ref{sec:annex-tutorial} introduce transformers y atención para lectores nuevos en el área.
    \item \textbf{Contrato de configuración}: la Sección~\ref{sec:config-centric-architecture} describe el esquema y la Sección~\ref{sec:schema-validation} sus reglas de validación.
    \item \textbf{Selección de arquitectura}: la Sección~\ref{sec:attention-families} compara las familias de mezcladores de secuencia; los Apéndices~\ref{sec:annex-transformer}, \ref{sec:annex-latent}, \ref{sec:annex-sparse} y~\ref{sec:annex-gated} recopilan la literatura de soporte, con los detalles de normalización, BitNet y memoria condicional en el Apéndice~\ref{sec:annex-norm-bitnet-mod} y los residuales con restricción de variedad en el Apéndice~\ref{sec:annex-mhc}.
    \item \textbf{Optimización}: la Sección~\ref{sec:optimizer-families} resume el enrutamiento de optimizadores; el Apéndice~\ref{sec:annex-optimizers} ofrece el análisis completo de las familias.
    \item \textbf{Despliegue e inferencia}: las Secciones~\ref{sec:quantization} y~\ref{sec:sbert} cubren el despliegue cuantizado y los embeddings de oraciones.
\end{enumerate}
